% !PRINTED 486
\weberpage{486}
%<<CONT
eine algebraische Gleichung besteht, und diese Gleichungen sind dort
für $m=7,\ 23,\ 31,\ 47,\ 71$ aufgestellt.

Setzen wir darin:
\[
\omega=\frac{-1}{\sqrt{-m}},\quad m\omega=\sqrt{-m},\quad f(\sqrt{-m})=\sqrt2\,x,
\]
so wird nach \S\,34:
\[
A=x^2+\frac{(-1)^{\frac{m+1}{8}}}{x},\quad B=4\,x+\frac{(-1)^{\frac{m+1}{8}}}{x}.
\]
Daraus ergibt sich z. B. für $m=23$, wo $A=1$ ist, die Gleichung:
\weq{1}{f(\sqrt{-23})=\sqrt2\,x,\quad x^3-x-1=0}
und für $m=31$:
\[
x^9-4\,x^6+3\,x^3-1=0.
\]
Dies ist zunächst eine kubische Gleichung für $x^3$; man spaltet daraus
aber leicht die kubische Gleichung für $x$ selbst ab:
\weq{2}{f(\sqrt{-31})=\sqrt2\,x,\quad x^3-x^2-1=0.}
Ähnlich ergeben sich die Gleichungen:
\weq{3}{f(\sqrt{-47})=\sqrt2\,x,\quad x^5-x^3-2\,x^2-2\,x-1=0,}
\weq{4}{f(\sqrt{-71})=\sqrt2\,x,\quad x^7-2\,x^6-x^5+x^4+x^3+x^2-x-1=0.}
Im letzten Falle, $m=71$, ist von der unmittelbar erhaltenen Gleichung
9ten Grades der der Aufgabe fremde Faktor $(x+1)^2$ abgesondert.

2. Um zu einer weiteren Berechnungsart zu gelangen, wenden wir die
Transformation zweiter Ordnung an. Wir haben zunächst allgemein
[für ein veränderliches $\omega$, \S\,34, (17)]:
\[
\begin{aligned}
f(2\omega)^8+f_2(2\omega)^8&=2\,\frac{f(\omega)^8+f_1(\omega)^8}{f_2(\omega)^4},\\
f(2\omega)^8-f_2(2\omega)^8&=\frac{16}{f_2(\omega)^8},
\end{aligned}
\]
woraus:
\[
\begin{aligned}
f(2\omega)^8&=\frac{f(\omega)^8+f_1(\omega)^8}{f_2(\omega)^4}+\frac{8}{f_2(\omega)^8},\\
f_2(2\omega)^8&=\frac{f(\omega)^8+f_1(\omega)^8}{f_2(\omega)^4}-\frac{8}{f_2(\omega)^8};
\end{aligned}
\]
%CONT>>
